\clearpage
\hypertarget{native-483-parthenope}{}
\begin{center}
\rule{0.55\linewidth}{0.4pt}\\[2pt]
Astronomische Nachrichten. Vol. XXX. No. 719. p. 374. 20 June 1850.\\[2pt]
\rule{0.55\linewidth}{0.4pt}
\end{center}
\begin{center}Observations and elements of Parthenope.\end{center}

Privy Councillor \textsc{Gauss} observed Parthenope with the meridian circle at the Göttingen Observatory,

\begin{center}
\setlength{\tabcolsep}{4pt}
\begin{tabular}{l|r|r|r}
 & Mean time & R.A. & Southern decl. \\\hline
May 31 & $10^{\rm u}\ 27^{\rm m}\ 38^{\rm s}$ & $225^\circ\ 52'\ 38''4$ & $9^\circ\ 48'\ 3''7$ \\
\end{tabular}
\end{center}

The air was somewhat hazy in the layers nearer the horizon, and with the wires at their faintest illumination the crossings were hard to make out.

\begin{center}
\setlength{\tabcolsep}{4pt}
\begin{tabular}{l|r|r|r}
June 2 & $10^{\rm u}\ 18^{\rm m}\ 1^{\rm s}6$ & $225^\circ\ 30'\ 3''6$ & $9^\circ\ 45'\ 59''8$ \\
\end{tabular}
\end{center}

